\documentclass[ngerman,xcolor={table,dvipsnames},smaller,compress,hyperref={bookmarks,colorlinks},handout]{beamer}%

\usepackage{url}
\usepackage{listings}
\usepackage[latin9]{inputenc}
\usepackage{textcomp}
%\usepackage{auto-pst-pdf}
\usepackage{graphicx}
\usepackage{pstricks}
\usepackage{pst-node}
\usepackage{pst-uml}
\usepackage{pst-tree}
\usepackage{pst-blur}

\usepackage{amsfonts}
\usepackage{amssymb}
\usepackage{amsmath}

% for multipage tables (xtab or longtable
\usepackage{longtable}
\usepackage{lscape}
\usepackage{booktabs}
\usepackage[tableposition=top]{caption}

\title{Moo in practice - System::Image::Update}
\author{Jens Rehsack}
\institute[Niederrhein.PM]{Niederrhein Perl Mongers}
\date{2015}

\usetheme[secheader]{Boadilla}
\setbeamertemplate{navigation symbols}{}

\newcommand{\perlfilename}[1]{{\color{cyan}{\textit{\begingroup \urlstyle{sf}\Url{#1}}}}}
\newcommand{\xsfilename}[1]{{\color{olive}{\textit{\begingroup \urlstyle{sf}\Url{#1}}}}}
\newcommand{\hfilename}[1]{{\color{olive}{\textit{\begingroup \urlstyle{sf}\Url{#1}}}}}
\newcommand{\cfilename}[1]{{\color{magenta}{\textit{\begingroup \urlstyle{sf}\Url{#1}}}}}

\makeatletter
\makeatother

\begin{document}

\psset{angleA=90,angleB=-90}
\lstset{language=Perl,
        basicstyle=\ttfamily,
        keywordstyle=\color{Maroon},
        commentstyle=\color{Blue}, 
        stringstyle=\color{Green},
        showstringspaces=false}

\AtBeginPart{\begin{frame}<beamer> \frametitle{Overview} \partpage \tableofcontents[current] \end{frame}}

\frame{\maketitle}

\part{Introduction}

\section{Introduction}

\subsection{Motivation}

\begin{frame}[t,fragile]{Motivation}

\begin{block}<1->{Moo and System::Image::Update}
\begin{itemize}
\item real world examples over far-fetched conceptuals
\item MooX::ConfigFromFile and MooX::Options provide way more features and flexibility than MooseX competitors
\item $ 2^{nd} $ generation of modern OO in Perl5
\end{itemize}
\end{block}

\begin{block}<2->{System::Image::Update}
\begin{itemize}
\item Provides an out-of-the-box solution for managing updates on embedded devices
\item Easy to re-use in several layers of the firmware
    \begin{itemize}
    \item self-sustaining (automatic) update management including forced updates (mind heartbleed)
    \item embeddable into middleware
    \item ability for shortcuts
    \end{itemize}
\item self-healing capabilities
\end{itemize}
\end{block}

\end{frame}

\subsection{Audience}

\begin{frame}[t,fragile]{Audience}
\begin{block}<1->{Audience}
\begin{itemize}
\item Developer who want to create or improve Perl5 software
\item Developer who want to learn how to develop modern OO with Perl5
\item DevOps or Integrators who are interested in embedded or automatic update concepts
\end{itemize}
\end{block}

\begin{block}<2->{Prerequisites of the Audience}
Following knowledge is expected:
\begin{itemize}
\item General knowledge about object oriented programming or concepts like
    \small
    \begin{itemize}
    \item classes
    \item objects
    \item polymorphism, inheritance and/or roles
    \item methods, class functions
    \item attributes, properties
    \end{itemize}
    \normalsize
\item slightly above basic Perl experience
\uncover<3->{\item ever heard of Smalltalk and its OO-concept is a strong bonus}
\end{itemize}
\end{block}
\end{frame}

\part{Moo basics}

\section{Modules}

\subsection{Classes and Roles}

\begin{frame}[t,fragile]

\begin{block}<1->{Classes in Moo}
\begin{itemize}
\item classes can be instantiated
\uncover<2->{\item one can inherit from classes}
\uncover<3->{\item one can aggregate classes}
\uncover<4->{\item distinguish naming conflicts is up to developer}
\end{itemize}
\end{block}

\begin{block}<5->{}
\scriptsize
\begin{lstlisting}[language=Perl,inputencoding=latin9,escapeinside={(*@}{@*)}]
{
  package System::Image::Update;

  use Moo;
  use MooX::Options with_config_from_file => 1;
  use IO::Async ();
  ...
  use File::ConfigDir::System::Image::Update qw(system_image_update_dir);
  around BUILDARGS => sub {...};
  sub run {..}
  sub collect_savable_config {}
  sub reset_config {}
  sub save_config {}
}
System::Image::Update->new_with_options->run;
\end{lstlisting}
\end{block}

\end{frame}

\begin{frame}[t,fragile]

\begin{block}<1->{Roles in Moo}
\begin{itemize}
\item roles describe a dedicated behavior (e.g. \texttt{logger})
\uncover<2->{\item roles can be composed into classes}
\uncover<3->{\item one can't inherit from roles - only consume}
\uncover<4->{\item roles cannot exist stand-alone}
\uncover<5->{\item roles are consumed once}
\uncover<6->{\item naming conflicts cause compile time error}
\end{itemize}
\end{block}

\begin{block}<7->{}
\scriptsize
\begin{lstlisting}[language=Perl,inputencoding=latin9,escapeinside={(*@}{@*)}]
{   package System::Image::Update::Role::HTTP;

    use Moo::Role; # now it's a role - no 'is a' relationship anymore

    sub do_http_request { ... }
    around collect_savable_config => sub {...};
}
{   package System::Image::Update::Role::Scan;

    use Moo::Role;
    with "System::Image::Update::Role::HTTP"; # consumes a role

    sub scan { my $self = shift; $self->do_http_request(...) };
}
\end{lstlisting}
\end{block}

\end{frame}

\section{Attributes}

\subsection{Attributes in Moo}

\begin{frame}[t,fragile]

\begin{block}<1->{Attributes in Moo - Overview}
\scriptsize
\begin{lstlisting}[language=Perl,inputencoding=latin9,escapeinside={(*@}{@*)}]
package System::Image::Update::Role::Scan;

use Moo::Role;

has(*@\pnode(0,0){A}{}@*) scan_(*@\pnode(0,0){B}{}@*)interval => ( is => "(*@\pnode(0,0){D}{}@*)ro", (*@\pnode(0,0){F}{}@*)default => 6*60*60 );
has update_mani(*@\pnode(0,0){C}{}@*)fest_uri => ( is => "(*@\pnode(0,0){E}{}@*)lazy(*@\pnode(0,0){G}{}@*)" );

1;
\end{lstlisting}
\end{block}

\begin{itemize}
\uncover<2->{\item use \rnode{a}{"has"} keyword to define a attribute}
\uncover<3->{\item attributes \rnode{b}{"scan\_interval"} and \rnode{c}{"update\_manifest\_uri"}}
\uncover<4->{\item those attributes are \rnode{d}{immutable}}
\uncover<5->{\item scan\_interval is initialized with a \rnode{f}{constant}}
\uncover<6->{\item update\_manifest\_uri is initialized on first access by a \rnode{g}{builder}}
\uncover<2>{\nccurve[linecolor=olive]{->}{a}{A}}
\uncover<3>{\nccurve[linecolor=teal]{->}{b}{B}\nccurve[linecolor=teal]{->}{c}{C}}
\uncover<4>{\nccurve[linecolor=blue]{->}{d}{D}\nccurve[linecolor=blue]{->}{d}{E}}
\uncover<5>{\nccurve[linecolor=cyan]{->}{f}{F}}
\uncover<6>{\nccurve[linecolor=cyan]{->}{g}{G}}
\end{itemize}

\end{frame}

\subsection{Attribute Options}

\begin{frame}[t,fragile]

\begin{block}<1->{Attribute options - Selection I}
\begin{description}
\item[is] \textbf{mandatory} behavior description
    \begin{description}
    \uncover<2->{\item[ro] defines the attribute is read-only}
    \uncover<3->{\item[rw] defined the attribute is read/writable}
    \uncover<4->{\item[lazy] defines the attribute is read-only with a lazy initialization, implies \texttt{builder => 1}}
    \end{description}
\uncover<5->{\item[required] when set to a true value, attribute must be passed on instantiation}
\uncover<6->{\item[isa] defines a subroutine (\texttt{coderef}) which is called to validate values to set}
\uncover<7->{\item[coerce] defines a subroutine (\texttt{coderef}) which forces attribute values}
\uncover<8->{\item[trigger] takes a subroutine (\texttt{coderef}) which is called anytime the attribute is set \\
    \uncover<9->{\textbf{special}: the value of $ 1 $ means to generate a (\texttt{coderef}) which calles the method \texttt{\_trigger\_\$\{attr\_name\}} (This is called \textit{attribute shortcut})}}
\uncover<10->{\item[default] subroutine (\texttt{coderef}) which is called to initialize an attribute}
\end{description}
\end{block}

\end{frame}

\begin{frame}[t,fragile]

\begin{block}<1->{Attribute options - Selection II}
\begin{description}[style=sameline]
\item[builder] takes a method name (\texttt{string}) which is called to initialize an attribute (supports \textit{attribute shortcut})
\uncover<2->{\item[init\_arg] Takes the name of the key to look for at instantiation time of the object. A common use of this is to make an underscored attribute have a non-underscored initialization name. undef means that passing the value in on instantiation is ignored.}
\uncover<3->{\item[clearer] takes a method name (\texttt{string}) which will clear the attribute (supports \textit{attribute shortcut})}
\uncover<4->{\item[predicate] takes a method name (\texttt{string}) which will return true if an attribute has a value (supports \textit{attribute shortcut})}
\end{description}
\end{block}

\end{frame}

\section{Methods}

\subsection{Method Examples}

\begin{frame}[t,fragile]

%\begin{columns}[t]
%\begin{column}{.8\textwidth}
\begin{block}{Methods in Moo}
\scriptsize
\begin{lstlisting}[language=Perl,inputencoding=latin9,escapeinside={(*@}{@*)}]
package System::Image::Update::Role::Async;

use IO::Async; use IO::Async::Loop;
use IO::Async::Timer::Absolute; use IO::Async::Timer::Countdown;

use Moo::Role;

has loop => ( is => "lazy", predicate => 1 );
sub _build_loop { return IO::Async::Loop->new() }

sub wakeup_at { my ( $self, $when, $cb_method ) = @_;
    my $timer;
    $self->loop->add($timer = IO::Async::Timer::Absolute->new(
        time      => $when,
        on_expire => sub { $self->$cb_method },
    ));
    $timer;
}
\end{lstlisting}
\end{block}

\onslide<2->{
\begin{itemize}
\item nothing like \texttt{MooseX::Declare} - pure Perl5 keywords are enough for plain methods
\end{itemize}
}

\end{frame}

\subsection{Method Modifiers}

\begin{frame}[t,fragile]

\begin{block}<1->{Method Modifiers}
Method modifiers are a convenient feature from the CLOS (Common Lisp Object System) world:
\begin{description}
\uncover<2->{\item[before] \texttt{before method(s) => sub \{ \ldots \}} \\
before is called before the method it is modifying. Its return value is totally ignored.}
\uncover<3->{\item[after] \texttt{after method(s) => sub \{ \ldots \}} \\
after is called after the method it is modifying. Its return value is totally ignored.}
\uncover<4->{\item[around] \texttt{around method(s) => sub \{ \ldots \}} \\
around is called instead of the method it is modifying. The method you're overriding is passed as \texttt{coderef} in the first argument.}
\uncover<5->{\item[-] No support for \texttt{su\pnode(0,0){A}{p}er}, \texttt{over\pnode(0,0){B}{}ride}, \texttt{in\pnode(0,0){C}{}ner} or \texttt{aug\pnode(0,0){D}{m}ent}}
\end{description}
\end{block}

\uncover<6->{In doubt \rnode{a}{\texttt{MooX::Override}} provides \texttt{override} and \texttt{super}}
\uncover<7->{and \rnode{b}{\texttt{MooX::Augment}} provides \texttt{augment} and \texttt{inner}}

\uncover<6>{\nccurve[linecolor=teal]{->}{a}{A}\nccurve[linecolor=teal]{->}{a}{B}}
\uncover<7>{\nccurve[linecolor=cyan]{->}{b}{C}\nccurve[linecolor=cyan]{->}{b}{D}}

\end{frame}

\begin{frame}[t,fragile]

\begin{block}<1->{Method Modifiers - Advantages}
\begin{itemize}
\item supersedes \texttt{\$self->SUPER::foo(@\_)} syntax
\uncover<2->{\item cleaner interface than SUPER}
\uncover<3->{\item allows multiple modifiers in single namespace}
\uncover<4->{\item also possible from within roles and not restricted to inheritance}
\uncover<5->{\item ensures that inherited methods invocation happens right (mostly - remember around)}
\uncover<6->{\item no need to change packages}
\end{itemize}
\end{block}

\end{frame}

\begin{frame}[t,fragile]

\begin{block}<1->{Methods Modifiers - around avoid calling \texttt{\$orig}}
\scriptsize
\begin{lstlisting}[language=Perl,inputencoding=latin9,escapeinside={(*@}{@*)}]
package Update::Status;

use strict; use warnings; use Moo;
extends "System::Image::Update";
around _build_config_prefix => sub { "sysimg_update" };
\end{lstlisting}
\end{block}

\begin{itemize}
\uncover<2->{\item captures control}
\uncover<3->{\item receives responsibility}
\uncover<4->{\item runtime of modified method completely eliminated}
\end{itemize}

\end{frame}

\begin{frame}[t,fragile]

\begin{block}<1->{Methods Modifiers - around modifying \texttt{\$orig} return value}
\scriptsize
\begin{lstlisting}[language=Perl,inputencoding=latin9,escapeinside={(*@}{@*)}]
package System::Image::Update::Role::Scan;

use strict; use warnings; use Moo::Role;

around collect_savable_config => sub {
    my $next                   = shift;
    my $self                   = shift;
    my $collect_savable_config = $self->$next(@_);

    $self->update_server eq $default_update_server
      or $collect_savable_config->{update_server} = $self->update_server;
    ...
    $collect_savable_config
}
\end{lstlisting}
\end{block}

\begin{itemize}
\uncover<2->{\item modifies only required part}
\uncover<3->{\item leaves most responsibility in modified method}
\uncover<4->{\item runtime of modified method added to this method's runtime}
\end{itemize}

\end{frame}

\part{System::Image::Update}

\section{Domain}

\subsection{Company}

\begin{frame}[t,fragile]

\begin{block}<1->{Rademacher Elektronik GmbH, Rhede}
  \begin{columns}
    \column[T]{.50\textwidth}
    \raisebox{\dimexpr.5\ht\strutbox-\totalheight-\partopsep-\topsep\relax}{%
      \includegraphics[width=5cm]{RDM-Factory-1-Front.eps}
    }%
    \column[T]{.50\textwidth}
    \uncover<2->{
      \begin{itemize}
      \item Successor of Guruplug based Homepilot 1
      \uncover<3->{\item Full update abilities (including boot loader, kernel and system components)}
      \uncover<4->{\item Multimedia features (Mediaplayer with HDMI video and Coaxial audio)}
      \uncover<5->{\item Closer Adoption to Rademacher Way: Control from Hardware development until Customer Service}
      \end{itemize}
    }
  \end{columns} 
\end{block}

\end{frame}

\subsection{Application}

\begin{frame}[t,fragile]

\begin{block}<1->{HomePilot 2}
  \begin{columns}
    \column[T]{.50\textwidth}
    \raisebox{\dimexpr.5\ht\strutbox-\totalheight-\partopsep-\topsep\relax}{%
      \includegraphics{Homepilot-angular.eps}
    }%
    \column[T]{.50\textwidth}
      \uncover<2->{
      \begin{itemize}
      \item service console moved from PHP to Perl5, PSGI and Dancer
      \uncover<3->{\item system management and automation full Perl5 powered}
      \uncover<4->{\item company infrastructure improved by Perl (eg. production of HP2)}
      \uncover<5->{\item created Yocto CPAN Layer for cross compiling lot's of CPAN modules}
      \uncover<6->{\item PoC use Perl6 + NativeCall to eliminate wrapper processes}
      \end{itemize}
    }
  \end{columns} 
\end{block}

\end{frame}

\section{Update Service}

\subsection{Self-Sustaining Daemon}

\begin{frame}[t,fragile]

\begin{block}<1->{State-Machine with togglable states and protected states}

\small
\resizebox{\textwidth}{!}{
  \psset{shadowcolor=black!70,blur=true}%
  \begin{pspicture}(25,18)%\psgrid
    \rput(4,17){\rnode{A}{\psovalbox[fillstyle=solid,fillcolor=yellow,shadow=true]{Begin}}}
    \uncover<2->{
        \rput(4,16){\rnode{B}{\psframebox[shadow=true]{Initialisations}}}
    }
    \uncover<3->{
        \rput(4,14){\rnode{bC}{}}
        \rput(4,13){\rnode{C}{\psframebox[shadow=true]{Scan}}}
            \uncover<4->{
                \rput(9,13){\rnode{D}{\psdiabox[fillstyle=solid,shadow=true]{Newer Image Information?}}}
            }
            \uncover<5->{
                \rput(15,13){\rnode{E}{\psframebox[shadow=true]{scan\_image\_info}}}
            }
            \uncover<6->{
                \rput(21,13){\rnode{F}{\psdiabox[fillstyle=solid,shadow=true]{Contain Update?}}}
            }
        \uncover<7->{
            \rput(1.5,15){\rnode{R}{\psframebox[shadow=true]{Reset}}}
        }
    }
    \uncover<8->{
        \rput(4,12){\rnode{bG}{}}
        \rput(4,11){\rnode{G}{\psframebox[shadow=true]{analyse\_update\_details}}}
    }
    \uncover<9->{
        \rput(4,9){\rnode{H}{\psdiabox[fillstyle=solid,shadow=true]{Want download?}}}
            \rput(9,9){\rnode{I}{\psframebox[shadow=true]{download}}}
            \rput(15,9){\rnode{J}{\psdiabox[fillstyle=solid,shadow=true]{Download successful?}}}
    }
    \uncover<10->{
        \rput(4,8){\rnode{bK}{}}
        \rput(4,7){\rnode{K}{\psdiabox[fillstyle=solid,shadow=true]{Verify download?}}}
            \rput(9,7){\rnode{L}{\psframebox[shadow=true]{verified}}}
    }
    \uncover<11->{
        \rput(4,6){\rnode{bM}{}}
        \rput(4,5){\rnode{M}{\psdiabox[fillstyle=solid,shadow=true]{Want apply?}}}
            \rput(9,5){\rnode{N}{\psframebox[shadow=true]{apply}}}
            \rput(15,5){\rnode{O}{\psovalbox[fillstyle=solid,fillcolor=orange,shadow=true]{Reboot}}}
    }


    % Links
    \uncover<2->{\ncline[linewidth=0.5pt,linestyle=solid]{->}{A}{B}}
    \uncover<3->{\ncline[linewidth=0.5pt,linestyle=solid]{->}{B}{C}}
    \uncover<4->{\ncline[linewidth=0.5pt,linestyle=solid]{->}{C}{D}}
    \uncover<5->{\ncline[linewidth=0.5pt,linestyle=solid]{->}{D}{E}\naput*[Ynodesep=0.5]{\small Yes}}
    \uncover<6->{\ncline[linewidth=0.5pt,linestyle=solid]{->}{E}{F}}
    \uncover<7->{
        \ncdiagg[angleA=90,lineAngle=0.01]{->}{D}{R}\naput*[Ynodesep=0.5]{\small \textcolor{red}{No}}
        \ncdiagg[angleA=90,lineAngle=0.01]{->}{F}{R}\naput*[Ynodesep=0.5]{\small \textcolor{red}{No}}
        \ncdiagg[angleA=-90,lineAngle=0.01]{->}{R}{bC}\naput*[Ynodesep=0.5]{\small scan}
    }
    \uncover<8->{
        \ncdiagg[angleA=-90,lineAngle=0.01]{->}{F}{bG}\naput*[Ynodesep=0.5]{\small Yes}
        \ncdiagg[angleA=-90,lineAngle=0.01]{->}{D}{bG}\naput*[Ynodesep=0.5,npos=0.50]{\small \textcolor{green}{But already}}
        \ncline[linewidth=0.5pt,linestyle=solid]{->}{bG}{G}
    }
    \uncover<9->{
        \ncline[linewidth=0.5pt,linestyle=solid]{->}{G}{H}
        \ncline[linewidth=0.5pt,linestyle=solid]{->}{H}{I}\naput*[Ynodesep=0.5]{\small Yes}
        \ncdiagg[angleA=180,lineAngle=-130,linestyle=solid]{->}{H}{bC}\nbput*[Ynodesep=0.5,npos=1.0]{\small \textcolor{green}{Not yet}}
        \ncline[linewidth=0.5pt,linestyle=solid]{->}{I}{J}
        \ncdiagg[angleA=-90,lineAngle=0.01]{->}{R}{bG}\nbput*[Ynodesep=0.5]{\small check}
    }
    \uncover<10->{
        \ncdiagg[angleA=-90,lineAngle=0.01]{->}{J}{bK}\naput*[Ynodesep=0.5]{\small Yes}
        \ncbar[angleA=0,angleB=0,lineAngle=0.01]{->}{G}{bK}\naput*[Ynodesep=0.25]{\small \textcolor{green}{But already}}
        \ncbar[angleA=0,angleB=0,linewidth=0.5pt,linestyle=solid]{->}{J}{R}\naput*[Ynodesep=0.5,npos=0.75]{\small \textcolor{red}{No}}
        \ncline[linewidth=0.5pt,linestyle=solid]{->}{bK}{K}
        \ncline[linewidth=0.5pt,linestyle=solid]{->}{K}{L}
        \ncbar[angleA=-180,angleB=-90,linewidth=0.5pt,linestyle=solid]{->}{K}{R}\naput*[Ynodesep=0.5]{\small \textcolor{red}{No}}
    }
    \uncover<11->{
        \ncdiagg[angleA=-90,lineAngle=0.01]{->}{L}{bM}
        \ncbar[angleA=180,angleB=180,lineAngle=0.01]{->}{G}{bM}\naput*[Ynodesep=0.5,npos=1.25]{\small \textcolor{green}{But already}}
        \ncline[linewidth=0.5pt,linestyle=solid]{->}{bM}{M}
        \ncline[linewidth=0.5pt,linestyle=solid]{->}{M}{N}
        \ncdiagg[angleA=180,lineAngle=-115,linestyle=solid]{->}{M}{bC}\nbput*[Ynodesep=0.5,npos=1.05]{\small \textcolor{green}{Not yet}}
        \ncline[linewidth=0.5pt,linestyle=solid]{->}{N}{O}
    }
    \end{pspicture}%
}%resizebox
\end{block}

\end{frame}

\begin{frame}[t,fragile]

\begin{block}<1->{State Control}
\scriptsize
\begin{lstlisting}[language=Perl,inputencoding=latin9,escapeinside={(*@}{@*)}]
package System::Image::Update;

use strict; use warnings; use Moo;
with "System::Image::Update::Role::Scan", "System::Image::Update::Role::Check", ...

has status => ( is => "rw", lazy => 1, builder => 1, predicate => 1,
    isa => sub { __PACKAGE__->can( $_[0] ) or die "Invalid status: $_[0]" }
);

sub _build_status { -f $_[0]->update_manifest ? "check" :
    $_[0]->has_recent_update && -e $_[0]->download_image ? "prove" : "scan";
}
\end{lstlisting}
\end{block}

\begin{itemize}
\uncover<2->{\item automatic recovering after down-state (power outage, Vodka party, \ldots)}
\uncover<3->{\item room for improvements like continue aborted download}
\uncover<4->{\item no direct path to ''download'' or ''apply'' to avoid mistakes}
\end{itemize}

\end{frame}

\begin{frame}[t,fragile]

\begin{block}<1->{State Control II}
\scriptsize
\begin{lstlisting}[language=Perl,inputencoding=latin9,escapeinside={(*@}{@*)}]
package System::Image::Update;

use strict; use warnings; use Moo;
with "System::Image::Update::Role::Scan",
     "System::Image::Update::Role::Check", ...

has status => ( ... );

around BUILDARGS => sub {
    my $next = shift; my $class = shift; my $params = $class->$next(@_);

    $params->{status} and $params->{status} eq "apply"
      and $params->{status} = "prove";
    $params->{status} and $params->{status} eq "prove"
      and $params->{recent_update}
      and $params->{recent_update}->{apply} = DateTime->now->epoch;
    $params;
};
\end{lstlisting}
\end{block}

\begin{itemize}
\uncover<2->{\item toggleable are ''download'' and ''prove''}
\uncover<3->{\item ''apply'' is protected by ''prove'' to ensure no corrupted image is applied}
\uncover<4->{\item protection needs to be improved before releasing to wildlife}
\end{itemize}

\end{frame}

\section{Middleware Interception}

\subsection{''ReadOnly''}

\begin{frame}[t,fragile]

\begin{block}<1->{Middleware Information Center}
\scriptsize
\begin{lstlisting}[language=Perl,inputencoding=latin9,escapeinside={(*@}{@*)}]
package Update::Status;

use strict; use warnings; use Moo;
extends (*@\pnode(0,0){A}{}@*)"System::Image::Update";

around _trigger(*@\pnode(0,0){C}{}@*)_recent_update => sub {};

around BUILDARGS => sub {
    my $next = shift; my $class = shift; my $params = $class->$next(@_);
    exists $params->{status} and (*@\pnode(0,0){B}{}@*)delete $params->{status};
    $params;
};
\end{lstlisting}
\end{block}

\begin{itemize}
\uncover<2->{\item \rnode{a}{derive from} ''System::Image::Update'' to get the real world picture}
\uncover<3->{\item \rnode{b}{ensure} no construction argument wastes the self-diagnostics}
\uncover<4->{\item \rnode{c}{prevent} triggering actions when recent update is found}
\end{itemize}
\uncover<2>{\nccurve[linecolor=olive]{->}{a}{A}}
\uncover<3>{\nccurve[linecolor=olive]{->}{b}{B}}
\uncover<4>{\nccurve[linecolor=olive]{->}{c}{C}}

\uncover<5>{\large What information?}

\end{frame}

\begin{frame}[t,fragile]

\begin{block}<1->{Middleware Information Center}
\scriptsize
\begin{lstlisting}[language=Perl,inputencoding=latin9,escapeinside={(*@}{@*)}]
package Update::Status;

use strict; use warnings; use Moo;
extends "System::Image::Update";

around _build(*@\pnode(0,0){A}{}@*)_config_prefix => sub { (*@\pnode(0,0){B}{}@*)"sysimg_update" };

around _trigger_recent_update => sub {};

around BUILDARGS => sub {
    my $next = shift; my $class = shift; my $params = $class->$next(@_);
    exists $params->{status} and delete $params->{status};
    $params;
};
\end{lstlisting}
\end{block}

\begin{itemize}
\uncover<2->{\item default builder guesses \rnode{a}{\texttt{''config\_prefix''}} from \texttt{\$0}}
\uncover<3->{\item \rnode{b}{override} with \texttt{\$0} from the daemon}
\end{itemize}
\uncover<2>{\nccurve[linecolor=teal]{->}{a}{A}}
\uncover<3>{\nccurve[linecolor=teal]{->}{b}{B}}

\end{frame}

\begin{frame}[t,fragile]

\begin{block}<1->{Middleware Delivery Center}
\scriptsize
\begin{lstlisting}[language=Perl,inputencoding=latin9,escapeinside={(*@}{@*)}]
package hp2sm;
use strict; use warnings; use Dancer2 ':syntax'; ...; use Update::Status;

get '/status' => sub {
    my $us = Update::Status->(*@\pnode(0,0){A}{}@*)new; my $status = (*@\pnode(0,0){B}{}@*)"idle";
    $us->has(*@\pnode(0,0){C}{}@*)_recent_update and $status = (*@\pnode(0,0){D}{}@*)"available";
    $us->stat(*@\pnode(0,0){E}{}@*)us eq "prove" and $status = (*@\pnode(0,0){F}{}@*)"downloading";
    $us->has_recent_update and -f $us->(*@\pnode(0,0){G}{}@*)download_image
      and $us->download_sums->{(*@\pnode(0,0){H}{}@*)size} == stat( $us->download_image )->(*@\pnode(0,0){I}{}@*)size
      and $status = (*@\pnode(0,0){J}{}@*)"proved";
    my $pl = (*@\pnode(0,0){K}{}@*)get_process_stats();
    firstval {$pl->process_name($_) =~ /(*@\pnode(0,0){L}{}@*)flash-device/} (0 .. $pl->entries()-1)
      and $status = "(*@\pnode(0,0){M}{}@*)applying";
    return $json->encode({status => $status});
};
\end{lstlisting}
\end{block}

\begin{itemize}
\uncover<2->{\item \rnode{a}{Load and guess} current status of ''System::Image::Update'' instance}
\uncover<3->{\item \rnode{b}{start with ''idle''} and \rnode{c1}{prove} from \rnode{d1}{earliest} to \rnode{m1}{latest}}
\uncover<7->{\item \rnode{j}{better} status overrides \rnode{f}{earlier} measures - \rnode{e}{\texttt{''prove''}} implies \rnode{c2}{\texttt{has\_recent\_update}}}
\uncover<10->{\item \rnode{l}{\texttt{''apply''}} is currently done by an external process}
\uncover<11->{\item use \rnode{k}{\texttt{Unix::Statgrab}} to grep for processes}
\end{itemize}

\uncover<2>{\nccurve[linecolor=teal]{->}{a}{A}}

\uncover<3-6>{\nccurve[linecolor=olive]{->}{b}{B}}
\uncover<4-6>{\nccurve[linecolor=olive]{->}{c1}{C}}
\uncover<5-6>{\nccurve[linecolor=RedViolet]{->}{d1}{D}}
\uncover<6-6>{\nccurve[linecolor=RedViolet]{->}{m1}{M}}

\uncover<7-9>{\nccurve[linecolor=olive]{->}{j}{J}}
\uncover<8-9>{\nccurve[linecolor=olive]{->}{f}{F}}
\uncover<9-9>{\nccurve[linecolor=cyan]{->}{e}{E}}
\uncover<9-9>{\nccurve[linecolor=cyan]{->}{c2}{C}}

\uncover<10>{\nccurve[linecolor=teal]{->}{l}{L}}
\uncover<11>{\nccurve[linecolor=teal]{->}{k}{K}}

\end{frame}

\subsection{''WriteBack''}

\begin{frame}[t,fragile]

\begin{block}<1->{Middleware Strikes Back}
\scriptsize
\begin{lstlisting}[language=Perl,inputencoding=latin9,escapeinside={(*@}{@*)}]
package hp2sm;
use strict; use warnings; use Dancer2 ':syntax'; ...; use Update::Status;

put '/status/downloading' => sub {
    my $us = Update::Status->new();
    $us->has_recent_update or return $json->encode( { result => "n/a" } );
    $us->status("download"); $us->save_config;
    system("svc -t /etc/daemontools/service/sysimg_update/");
    return $json->encode( { result => "ok" } );
};
\end{lstlisting}
\end{block}

\begin{block}<2->{Middleware Information Boosted Persistency}
\scriptsize
\begin{lstlisting}[language=Perl,inputencoding=latin9,escapeinside={(*@}{@*)}]
package Update::Status;
use strict; use warnings; use Moo;
extends "System::Image::Update";

around collect_savable_config => sub {
    my $next = shift; my $self = shift; my $save_cfg = $self->$next(@_);
    $self->has_status and $save_cfg->{status} = $self->status;
    $self->has_download_file
      and $save_cfg->{download_file} = $self->download_file;
    $save_cfg; };
\end{lstlisting}
\end{block}

\end{frame}

\section{Glue}

\subsection{Into Target Filesystem}

\begin{frame}[t,fragile]

\begin{block}<1->{system-image-update\_git.bb top}
\scriptsize
\begin{lstlisting}[language=bash,inputencoding=latin9,escapeinside={(*@}{@*)}]
DESCRI(*@\pnode(0,0){A}{}@*)PTION = "System::Image::Update helps managing updates of OS images ..."
SRC_URI(*@\pnode(0,0){B}{}@*) = "git://github.com/rehsack/System-Image-Update.git;rev=646fa928... \
           file://run file://sysimg_update.json"
RDEPENDS_${PN}(*@\pnode(0,0){C}{}@*) += "archive-peek-libarchive-perl crypt-ripemd160-perl"
RDEPENDS_${PN} += "datetime-format-strptime-perl"
RDEPENDS_${PN} += "log-any-adapter-dispatch-perl"
RDEPENDS_${PN} += "file-configdir-system-image-update-perl"
RDEPENDS_${PN}(*@\pnode(0,0){D}{}@*) += "moo-perl moox-configfromfile-perl moox-log-any-perl"
RDEPENDS_${PN} += "moox-options-perl net-async-http-perl"
RDEPENDS_${PN} += "digest-md5-perl digest-md6-perl"
RDEPENDS_${PN} += "digest-sha-perl digest-sha3-perl"
RDEPENDS_${PN}(*@\pnode(0,0){E}{}@*) += "daemontools system-image"
S = "${WORKDIR}/(*@\pnode(0,0){F}{}@*)git"
BBCLASSEXTEND = (*@\pnode(0,0){G}{}@*)"native"

inherit (*@\pnode(0,0){H}{}@*)cpan
do_configure_append() {
    oe_runmake (*@\pnode(0,0){I}{}@*)manifest
}
\end{lstlisting}
\end{block}

\begin{itemize}
\uncover<2->{\item typical  \rnode{a}{package} stuff \ldots, \uncover<3->{like \rnode{c}{runtime dependencies}}}
\uncover<4->{\item git checkouts need \rnode{f}{adoption} of source path}
\uncover<5->{\item build as \rnode{h}{any cpan} package is built, but allow \rnode{g}{native} packages and create missing \rnode{i}{\texttt{MANIFEST}}}
\end{itemize}

\uncover<2-3>{\nccurve[linecolor=teal]{->}{a}{A}}
\uncover<2-3>{\nccurve[linecolor=teal]{->}{a}{B}}
\uncover<3-3>{\nccurve[linecolor=teal]{->}{c}{C}}
\uncover<3-3>{\nccurve[linecolor=teal]{->}{c}{D}}
\uncover<3-3>{\nccurve[linecolor=teal]{->}{c}{E}}

\uncover<4-4>{\nccurve[linecolor=RedViolet]{->}{f}{F}}

\uncover<5-6>{\nccurve[linecolor=cyan]{->}{h}{H}}
\uncover<5-6>{\nccurve[linecolor=cyan]{->}{g}{G}}
\uncover<6-6>{\nccurve[linecolor=RedViolet]{->}{i}{I}}

\end{frame}

\begin{frame}[t,fragile]

\begin{block}<1->{system-image-update\_git.bb bottom}
\scriptsize
\begin{lstlisting}[language=bash,inputencoding=latin9,escapeinside={(*@}{@*)}]
SERVICE_ROOT = "${sysconfdir}/daemontools/service"
SYSUPDT_SERVICE_(*@\pnode(0,0){A}{}@*)DIR = "${SERVICE_ROOT}/sysimg_update"

do_install_append() {
    install -d -m 755 ${D}${sysconfdir}
    install -m 0644 ${WORKDIR}/(*@\pnode(0,0){C}{}@*)sysimg_update.json ${D}${sysconfdir}

    install -d ${D}${SYSUPDT_SERVICE_DIR}
    install -m 0755 ${WORKDIR}/(*@\pnode(0,0){B}{}@*)run ${D}${SYSUPDT_SERVICE_DIR}/run
}
FILES_${PN} (*@\pnode(0,0){D}{}@*)+= "${sysconfdir}"
\end{lstlisting}
\end{block}

\begin{itemize}
\uncover<2->{\item define \rnode{a}{location of startup} scripts \uncover<3->{and \rnode{b}{install to there}}}
\uncover<4->{\item \rnode{c}{install} configuration file}
\uncover<5->{\item tell bitbake to put files from \texttt{\$\{sysconfdir\}} \rnode{d}{into package}}
\end{itemize}

\uncover<2-3>{\nccurve[linecolor=teal]{->}{a}{A}}
\uncover<3-3>{\nccurve[linecolor=cyan]{->}{b}{B}}

\uncover<4-4>{\nccurve[linecolor=olive]{->}{c}{C}}
\uncover<5-5>{\nccurve[linecolor=olive]{->}{d}{D}}

\end{frame}

\section{Configuration file}

\subsection{Distributions and Moo Initialization Glue}

\begin{frame}[t,fragile]

\begin{block}<1->{sysimg\_update.json}
\scriptsize
\begin{lstlisting}[language=Java,inputencoding=latin9,escapeinside={(*@}{@*)}]
{
    "(*@\pnode(0,0){A}{}@*)log_adapter" : [
        "Dispatch",
        "outputs", [
            [ "File",   "min_level", "debug", "filename",
              "/var/log/sysimg_update.log", "newline", 1, "mode", ">>" ],
            [ "File",   "min_level", "error", "filename",
              "/var/log/sysimg_update.error", "newline", 1, "mode", ">>" ],
            [ "Screen", "min_level", "notice", "newline", 1, "stderr", 1 ]
        ]
    ],
    "(*@\pnode(0,0){B}{}@*)update_manifest_dirname" : "/rwmedia/update/",
    "(*@\pnode(0,0){C}{}@*)http_user" : "b01f..."
}
\end{lstlisting}
\end{block}

\begin{itemize}
\uncover<2->{\item Provides \rnode{a}{settings} for \texttt{Log::Any} (mind \texttt{\_trigger\_log\_adapter} in \texttt{System::Image::Update::Role::Logging} consuming \texttt{MooX::Log::Any})}
\uncover<3->{\item redirect place to \rnode{b}{store update} manifest (files)}
\uncover<4->{\item Provide \rnode{c}{authentication} to update server for development boxes (avoid \texttt{builder} being called)}
\end{itemize}

\uncover<2-2>{\nccurve[linecolor=olive]{->}{a}{A}}
\uncover<3-3>{\nccurve[linecolor=teal]{->}{b}{B}}
\uncover<4-4>{\nccurve[linecolor=RedViolet]{->}{c}{C}}

\end{frame}
\part{Finish}

\section{Conclusion}

\begin{frame}[t,fragile]

\begin{block}<1->{Conclusion}
\begin{itemize}
\item lazy attributes allow designing a multi-stage initialization phase
\uncover<2->{\item benefit of common runtime (faster load) when using}
\uncover<3->{\item improve design by}
\begin{itemize}
\uncover<4->{\item using roles for behavioral design (avoid duck typing)}
\uncover<5->{\item using explicit patterns for clear separation of concerns}
\uncover<6->{\item express intentions clearer for method overloading by using \textit{method modifiers}}
\end{itemize}
\end{itemize}
\end{block}

\end{frame}

\section{Resources}

\subsection{Resources}

% RESOURCES
\begin{frame}[t,fragile]{Resources}
\begin{block}<1->{Software on MetaCPAN}
\url{https://metacpan.org/pod/Moo} \\
\url{https://metacpan.org/search?q=MooX} \\
\url{https://metacpan.org/pod/MooX::Options} \\
\url{https://metacpan.org/pod/MooX::ConfigFromFile} \\
\url{https://metacpan.org/pod/IO::Async}
\end{block}

\begin{block}<2->{Software on GitHub}
\url{https://github.com/moose/Moo} \\
\url{https://github.com/rehsack/System-Image-Update} \\
\url{https://github.com/perl5-utils/File-ConfigDir-System-Image-Update}
\end{block}

\begin{block}<3->{Software for Cross-Building Perl-Modules}
\url{https://www.yoctoproject.org/} \\
\url{https://github.com/rehsack/meta-cpan}
\end{block}
\end{frame}

\begin{frame}[t,fragile]{Resources}
\begin{block}<1->{IRC}
\url{irc://irc.perl.org/#moose} \\
\url{irc://irc.perl.org/#web-simple} \\
\url{irc://irc.perl.org/#dancer} \\
\url{irc://irc.freenode.org/#yocto}
\end{block}

\begin{block}<2->{Hints}
\url{http://sck.pm/WVO} \# proper usage of the roles in perl \\
\url{https://metacpan.org/pod/Moo\#CLEANING-UP-IMPORTS}
\end{block}
\end{frame}

\subsection{Thank you}

% THANKS
\begin{frame}[fragile]{Thank You For Listening}
\begin{block}<1->{Questions?}
Jens Rehsack \textless{}\href{mailto:rehsack@cpan.org}{rehsack@cpan.org}\textgreater{}
\end{block}
\end{frame}

\end{document}
